\documentclass[12 pt, letterpaper]{hmcpset}
\usepackage{hw-preamble}

\name{Jayden Thadani}
\class{MATH 135 --- Complex Analysis}
\assignment{Homework 8}
\duedate{22nd November 2024}

\begin{document}

\begin{problem}[2.]
	Find the order of growth of the following entire functions:
	\begin{enumerate}
		\item $p(z)$ where $p$ is a polynomial.
		\item $e^{bz^{n}}$ for $b \neq 0$.
		\item $e^{e^{z}}$.
	\end{enumerate}
\end{problem}

\begin{solution}
	\begin{enumerate}
		\item
			\[
				\boxed{
					\rho = 0
				}
			\]

			Note that $e^{\lvert z \rvert^{\epsilon}}$ is a Taylor series $\sum_{n = 0}^{\infty} (\lvert z \rvert^{\varepsilon})^{n} / n!$. Since all the terms in the Taylor series are positive, we have
			\[
				\lvert z \rvert^{\varepsilon} + \frac{\lvert z \rvert^{2 \varepsilon}}{2!} + \dots + \frac{\lvert z \rvert^{m \varepsilon}}{m!} < e^{\lvert z \rvert^{\varepsilon}}
			\]
			for all $m$. Since large enough $m$ gives us $m \varepsilon > n$ for any $n$ by the Archimedean property, any polynomial $p$ of degree $n$ gives
			\[
				\lvert p(z) \rvert \le \lvert z \rvert^{\varepsilon} + \frac{\lvert z \rvert^{2 \varepsilon}}{2!} + \dots + \frac{\lvert z \rvert^{m \varepsilon}}{m!} \le e^{\lvert z \rvert^\varepsilon}
			\]
			for large enough $z$.
		\item
			\[
				\boxed{
					\rho = n	
				}
			\]

			We know that any $\lvert a \rvert \lvert z \rvert^{r}$ is eventually less than $\lvert b \rvert \lvert z \rvert^{n}$ for $r < n$. Thus, $\lvert e^{a z^{r}} \rvert$ is eventually less than $\lvert e^{b z^{n}} \rvert$. We also know that $\lvert e^{b z^{n}} \rvert \le e^{\lvert b \rvert \lvert z \rvert^{n}}$. Thus, $\rho \le n$ and $r \le \rho$ for any $r < n$. Thus, $\rho = n$.
		\item
			\[
				\boxed{
					\rho = \infty
				}
			\]

			We know that at large enough $z$, $\lvert e^{z} \rvert \ge b \lvert z \rvert^{r}$ for any $b, r$. Thus,
			\[
				e^{e^{z}} \ge e^{\lvert e^{z} \rvert} \not \le e^{b \lvert z \rvert^{r}} \ge A e^{B \lvert z \rvert^{r}}
			\]
			with the right choice of $b$. Thus there is no real number $\rho$ such that $\lvert e^{e^{z}} \rvert \le \lvert e^{B z^{\rho}} \rvert$.
	\end{enumerate}
\end{solution}

\pagebreak

\begin{problem}[3.]
	Show that if $\tau$ is fixed with $\operatorname{Im} \tau > 0$, then the Jacobi theta function
	\[
		\Theta(z \mid \tau) = \sum_{n = -\infty}^{n = \infty} e^{\pi i n^{2} \tau} e^{2 \pi i n z}
	\]
	is of order $\le 2$ as a function of $z$.
\end{problem}

\begin{solution}
	Write $\operatorname{Im} \tau = t > 0$ and $\operatorname{Im} z = y$. We can bound the theta function by the triangle inequality and proceed from there ---
	\begin{align*}
		\lvert \Theta(z \mid \tau) \rvert & = \left \lvert \sum_{n = - \infty}^{\infty} e^{\pi i n^{2} \tau} e^{2 \pi i n z} \right \rvert \\
						  & \le \sum_{n = -\infty}^{\infty} \lvert e^{\pi i n^{2} \tau} e^{2 \pi i n z} \rvert \\
						  & = \sum_{n = -\infty}^{\infty} \lvert e^{-\pi n^{2} t - 2 \pi n y} \rvert \\
						  & \le \sum_{n = -\infty}^{\infty} \lvert e^{\pi(- n^{2} t + 2 n \lvert z \rvert)} \rvert \\
						  & \le 2 \sum_{n = 0}^{\infty} \lvert e^{\pi(-n^{2} t + 2 n \lvert z \rvert)} \rvert.
	\end{align*}

	Here, we use the hint that $-n^{2} t + 2 n \lvert z \rvert \le - n^{2} t / 2$ when $n \ge 4 \lvert z \rvert / t$ and use this to split the sum at $n = \lfloor 4 \lvert z \rvert / t \rfloor$ (while this does depend on $\lvert z \rvert$, we will show that we can still bound the sum). Splitting the sum as described, we obtain
	\begin{align*}
		\lvert \Theta(z \mid \tau) \rvert & \le 2 \left ( \sum_{n = 0}^{\lfloor 4 \lvert z \rvert / t \rfloor} \lvert e^{\pi (-n^{2} t + 2 n \lvert z \rvert)} \rvert + \sum_{n = \lfloor 4 \lvert z \rvert / t \rfloor + 1}^{\infty} \lvert e^{\pi (-n^{2} t + 2 n \lvert z \rvert)} \rvert \right ) \\
						  & \le 2 \left ( \sum_{n = 0}^{\lfloor 4 \lvert z \rvert / t \rfloor} \lvert e^{\pi (-n^{2} t + 2 n \lvert z \rvert)} \rvert + \sum_{n = \lfloor 4 \lvert z \rvert / t \rfloor + 1}^{\infty} \lvert e^{- \pi n^{2} t/2} \rvert \right )
	\end{align*}
	Since the sum on the right is always less than a constant by comparison with the Gaussian integral. Thus we only need to show that the sum on the left is of order at most $2$.

	Note that in the sum up to $n = 4 \lvert z \rvert / t$, $e^{- \pi n^{2} t}$ is maximised when $n = 0$ and $e^{2 \pi n \lvert z \rvert}$ is maximised when $n = \lfloor 4 \lvert z \rvert / t \rfloor$. Thus, the summand $e^{- \pi n^{2} t} e^{2 \pi n \lvert z \rvert}$ is bounded by the product of its maximised factors. In particular
	\[
		e^{\pi(-n^{2} t + 2 n \lvert z \rvert)} \le e^{8 \pi n \lvert z \rvert^{2} / t}.
	\]
	Bounding the sum by the maximum summand, we get
	\[
		\lvert \Theta(z \mid \tau) \rvert \le \frac{4 \lvert z \rvert}{t} e^{8 \pi n \lvert z \rvert^{2} / t}.
	\]
	Since sufficiently large positive $A, B \in \mathbb{R}$ will give
	\[
		\frac{4 \lvert z \rvert}{t} e^{8 \pi n \lvert z \rvert / t} \le A e^{B \lvert z \rvert^{2 + \varepsilon}}
	\]
	for any $\varepsilon > 0$, $\Theta(z \mid \tau)$ must have growth of order less than or equal to $2$.
\end{solution}

\pagebreak

\begin{problem}[7.]
	Establish the following properties of infinite products
	\begin{enumerate}
		\item Show that if $\sum \lvert a_{n} \rvert^{2}$ converges, then the product $\prod (1 + a_{n})$ converges to a non-zero limit if and only if $\sum a_{n}$ converges.
		\item Find an example of a sequence of complex numbers $\{ a_{n} \}$ such that $\sum a_{n}$ converges but $\prod (1 + a_{n})$ diverges.
	\end{enumerate}
\end{problem}

\begin{solution}
	\begin{enumerate}
		\item
		\item By analysing the proof that absolutely convergent $\sum a_{n}$  implies convergent  $\prod (1 + a_{n})$ we see that we want convergent $\sum a_{n}$ such that  $\sum \log (1 + a_{n})$ diverges. Such a series is given by
			\[
				a_{n} = \frac{(-1)^{n}}{\sqrt{n}}.
			\]

			By the alternating series test $\sum a_{n}$ converges.

			If we write out $\log(1 + a_{n})$ using Taylor series we see that
			 \[
				\log (1 + a_{n}) = a_{n} - \frac{a_{n}^{2}}{2} + \dots = a_{n} - \frac{1}{2n} + \dots
			\]
			where the rest of the terms are $O(n^{-3/2})$. Since the higher order terms converge in series by the  $p$-test, $\sum a_{n}$ converges as as shown earlier, and $\sum -2/n$ diverges, the series must diverge. Thus
			\[
				\prod_{n = 1}^{\infty} (1 + a_{n}) = e^{\sum \log(1 + a_{n})}
			\]
			diverges.
	\end{enumerate}
\end{solution}

\pagebreak

\begin{problem}[8.]
	Prove that for every $z$ the product below converges, and
	\[
		\cos{(z / 2)} \cos{(z / 4)} \cos{(z / 8)} \dots = \prod_{k = 1}^{\infty} \cos{(z / 2^{k})} = \frac{\sin{z}}{z}
	\]
\end{problem}

\begin{solution}
	Consider $\sin{(z/ 2^{n})}$ times the partial product ---
	\[
		\sin{(z/2^{n})} \prod_{k = 1}^{n} \cos{(z/2^{k})}.
	\]
	Using the fact that
	\[
		\sin{(z / 2^{n})} \cos{z / 2^{n}} = \frac{\sin{(z / 2^{n - 1})}}{2}
	\]
	and proceeding by induction, we get
	\[
		\sin{(z / 2^{n})} \prod_{k = 1}^{n} \cos{(z / 2^{k})} = \frac{\sin{z}}{2^{n}}.
	\]

	With some algebraic manipulation, we get the partial product
	\[
		\prod_{k = 1}^{n} \cos{(z / 2^{k})} = \frac{\sin{z}}{z} \frac{z / 2^{n}}{\sin{(z / 2^{n})}}.
	\]
	Using the fact that $\lim_{z \to 0} \sin (z) / z = 1$, we get
	\[
		\lim_{n \to 0} \prod_{k = 1}^{n} \cos{(z / 2^{k})} = \lim_{n \to 0} \frac{\sin{z}}{z} \frac{z/2^{n}}{\sin{(z / 2^{n})}} = \frac{\sin{z}}{z}.
	\]
\end{solution}

\pagebreak

\begin{problem}[9.]
	Prove that if $\lvert z \rvert < 1$, then
	\[
		(1 + z)(1 + z^{2})(1 + z^{4})(1 + z^{8}) \dots = \prod_{k = 0}^{\infty} (1 + z^{2^{k}}) = \frac{1}{1 - z}.
	\]
\end{problem}

\begin{solution}
	Consider the partial product
	\[
		\prod_{k = 0}^{n} (1 + z^{2^{k}}) = \sum_{k = 0}^{2^{n + 1} - 1} z^{k}.
	\]
	Thus, the partial products form a subsequence of the partial sums of the geometric series. Since this sequence converges for $\lvert z \rvert < 1$, the subsequence also converges to the geometric series. That is,
	\[
		\prod_{k = 0}^{\infty} (1 + z^{2^{k}}) = \lim_{n \to \infty} \prod_{k = 0}^{n} (1 + z^{2^{k}}) = \lim_{n \to \infty} \sum_{k = 0}^{2^{n + 1} - 1} z^{k} = \frac{1}{1 - z}.
	\]
\end{solution}

\end{document}
